\documentclass[10pt]{article}
\usepackage{amsmath,amsfonts,amssymb,mathtools,mathrsfs,bm,bbm,array,xcolor,geometry}
\allowdisplaybreaks
\setlength{\parindent}{0pt}
\setlength{\parskip}{4pt}
\emergencystretch=4em

\geometry{paperwidth=215mm,paperheight=290mm,hmargin=15mm,vmargin=16mm}
\begin{document}
\section*{P0024 Food Basket Procurement and Delivery with Nutrition and Palatability}
\subsection*{Mathematical formulation}
\begin{equation}
\min_{\boldsymbol{x},\, y,\, \boldsymbol{F}}
\sum_{i \in \mathcal{N}_{\mathcal S}}
\sum_{j \in \mathcal{N}_{\mathcal T} \cup \mathcal{N}_{\mathcal D}}
\sum_{k \in \mathcal K} p_{ik}^P F_{ijk}
+
\sum_{i \in \mathcal{N}_{\mathcal S} \cup \mathcal{N}_{\mathcal T}}
\sum_{j \in \mathcal{N}_{\mathcal T} \cup \mathcal{N}_{\mathcal D}}
\sum_{k \in \mathcal K} p_{ijk}^T F_{ijk}
\tag{6a}
\end{equation}

\begin{equation}
\sum_{j \in \mathcal N_{\mathcal T}\cup\mathcal N_{\mathcal D}} F_{ijk}
=
\sum_{j \in \mathcal N_{\mathcal S}\cup\mathcal N_{\mathcal T}} F_{jik}
\quad \forall i \in \mathcal N_{\mathcal T},\ \forall k \in \mathcal K
\tag{6b}
\end{equation}

\begin{equation}
\sum_{j \in \mathcal N_{\mathcal S} \cup \mathcal N_{\mathcal T}} \gamma F_{jik}
= d_i x_k \, days
\quad \forall i \in \mathcal N_{\mathcal D},\ \forall k \in \mathcal K
\tag{6c}
\end{equation}

\begin{equation}
\sum_{k \in \mathcal K} Nutval_{kl} x_k \ge Nutreq_l
\quad \forall l \in \mathcal L
\tag{6d}
\end{equation}

\begin{equation}
x_{\mathrm{salt}} = 5
\tag{6e}
\end{equation}

\begin{equation}
x_{\mathrm{sugar}} = 20
\tag{6f}
\end{equation}

\begin{equation}
y \ge t
\tag{6g}
\end{equation}

\begin{equation}
y = \hat h(\boldsymbol{x})
\tag{6h}
\end{equation}

\begin{equation}
F_{ijk},\, x_k \ge 0
\quad \forall i,j \in \mathcal N,\ \forall k \in \mathcal K
\tag{6i}
\end{equation}

\begin{equation}
a_k^{(0)} = x_k
\quad \forall k \in \mathcal K
\tag{Emb-a}
\end{equation}

\begin{equation}
s_r^{(\ell)} =
\sum_{q \in \mathcal R_{\ell-1}} W_{rq}^{(\ell)} a_q^{(\ell-1)} + b_r^{(\ell)}
\quad \forall \ell=1,\dots,H,\ \forall r \in \mathcal R_\ell
\tag{Emb-b}
\end{equation}

\begin{equation}
a_r^{(\ell)} \ge s_r^{(\ell)}
\quad \forall \ell=1,\dots,H,\ \forall r \in \mathcal R_\ell
\tag{Emb-c}
\end{equation}

\begin{equation}
a_r^{(\ell)} \le s_r^{(\ell)} - L_r^{(\ell)} \left(1-\delta_r^{(\ell)}\right)
\quad \forall \ell=1,\dots,H,\ \forall r \in \mathcal R_\ell
\tag{Emb-d}
\end{equation}

\begin{equation}
a_r^{(\ell)} \le U_r^{(\ell)} \delta_r^{(\ell)}
\quad \forall \ell=1,\dots,H,\ \forall r \in \mathcal R_\ell
\tag{Emb-e}
\end{equation}

\begin{equation}
a_r^{(\ell)} \ge 0
\quad \forall \ell=1,\dots,H,\ \forall r \in \mathcal R_\ell
\tag{Emb-f}
\end{equation}

\begin{equation}
y =
\sum_{q \in \mathcal R_H} W_q^{(H+1)} a_q^{(H)} + b^{(H+1)}
\tag{Emb-g}
\end{equation}

\begin{equation}
\delta_r^{(\ell)} \in \{0,1\}
\quad \forall \ell=1,\dots,H,\ \forall r \in \mathcal R_\ell
\tag{Emb-h}
\end{equation}

\[
\sum_jF_{ijk}\le Cap^P_{ik}\ (i\in N_S),\quad
\sum_kF_{ijk}\le Cap^T_{ij},\quad0\le x_k\le U_k,\quad
L_g\le\sum_{k\in K_g}x_k\le U_g.
\]

\subsection*{Definitions and model scope}
A food-aid organization must supply the same daily food basket to demand locations and arrange its procurement and delivery at minimum combined procurement and transport cost. The network nodes $\mathcal N$ comprise sources $\mathcal N_{\mathcal S}$, transfer locations $\mathcal N_{\mathcal T}$ and demand locations $\mathcal N_{\mathcal D}$, with foods in $\mathcal K$ and nutrients in $\mathcal L$. Daily food quantities $x_k$ and directed shipment quantities $F_{ijk}$ may be any nonnegative real amounts. Procurement is charged on food leaving each source for transfer or demand locations at the source food price $p_{ik}^P$. Transport is charged at $p_{ijk}^T$ on shipments from sources or transfer locations to transfer or demand locations.

At each transfer location, total outflow of each food must equal its total inflow. At each demand location, incoming food multiplied by the conversion factor $\gamma$ must exactly equal the product of daily demand $d_i$, that food's daily basket quantity and the planning duration $days$. Total dispatch of a food from a source cannot exceed its procurement capacity $Cap^P_{ik}$, and total food tonnage on each transport arc cannot exceed $Cap^T_{ij}$.

For every nutrient, the sum of each food's nutrient content $Nutval_{kl}$ multiplied by its daily quantity must reach at least the requirement $Nutreq_l$. The daily salt quantity $x_{\mathrm{salt}}$ is fixed at five and the sugar quantity $x_{\mathrm{sugar}}$ at twenty, in the basket's stated units. Each food must also remain below its individual upper limit $U_k$. For every food category $K_g$, the combined daily quantities of its foods must lie between $L_g$ and $U_g$.

The basket's palatability score $y$ must reach the threshold $t$ and is determined by the already trained predictor $\hat h(\boldsymbol{x})$; procurement planning does not retrain it. Food quantities enter as $a_k^{(0)}$ and pass through $H$ hidden layers, each with processing units $\mathcal R_\ell$. A unit multiplies previous-layer outputs $a_q^{(\ell-1)}$ by trained weights $W_{rq}^{(\ell)}$, sums them and adds bias $b_r^{(\ell)}$ to obtain its raw response $s_r^{(\ell)}$. Its output $a_r^{(\ell)}$ is the larger of that response and zero. Raw responses remain within valid lower and upper bounds $L_r^{(\ell)}$ and $U_r^{(\ell)}$.

Each processing unit is either active or inactive, recorded by $\delta_r^{(\ell)}$ taking one or zero. Its output must be nonnegative and no smaller than the raw response, while not exceeding the raw response minus the lower bound multiplied by the inactive indicator, or the upper bound multiplied by the active indicator. The final score is the last-layer outputs weighted by $W_q^{(H+1)}$, plus output bias $b^{(H+1)}$. Foods, nutrition information, demand, costs, capacities and the trained predictor are fixed beforehand. If delivery, nutritional, quantity-range and palatability requirements cannot all be met, no implementable basket and delivery arrangement exists.

\subsection*{Inline notation and relations}
$\hat h(\boldsymbol{x})$; $x_k$; $k$; $F_{ijk}$; $y$; $\mathcal N_{\mathcal S}$; $\mathcal N_{\mathcal T}$; $\mathcal N_{\mathcal D}$; $\mathcal K$; $\mathcal L$; $p_{ik}^P$; $p_{ijk}^T$; $\gamma$; $d_i$; $days$; $Nutval_{kl}$; $l$; $Nutreq_l$; $t$; $\mathcal R_\ell$; $\ell=1,\dots,H$; $W_{rq}^{(\ell)}$; $b_r^{(\ell)}$; $L_r^{(\ell)} \le s_r^{(\ell)} \le U_r^{(\ell)}$; $Cap^P$; $Cap^T$; $U$; $K_g$; $L_g,U_g$.
\end{document}